%%==========================================================================
\section{Ordered Event-Based Reformulation for a Flight-hour Maintenance}
\label{sec:event_reformulation}
%%==========================================================================

In this paper, we present an ordered event-based model for
flight-hour maintenance. It removes day-indexed maintenance decision states
and does not introduce the
standalone successor-arc event formulation. Flights are processed in
nondecreasing departure-time order, and maintenance is represented by a binary
event attached to the flight that triggers it. Thus the design-note variables
$x_{ij}$ and $z_{ij}$ are represented below through an ordered position index
$r$ that maps to a flight index $i_r$.

\subsection{Ordered event representation}

For flight $i$, let $o_i$ and $a_i$ be its departure and arrival airports,
$p_i$ its flying time, and $\tD{i}$ and $\tI{i}$ its departure and arrival
timestamps (the same macros as in \Cref{sec:basic}). For aircraft $j$, let $a^0_j$ be its initial airport.

Define the ordered-position index set
\[
\mathcal{R}=\{1,\ldots,n_f\},
\]
and let $(i_1,\ldots,i_{n_f})$ be flights sorted by $(\tD{i},i)$. Position
$r\in\mathcal{R}$ refers to flight $i_r$. This is the only role of $r$: it is
an order index over flights. Define
\[
\mathcal{R}^M=\{r\in\mathcal{R}: i_r\in\F^M\},
\]
and let $\mathcal{Z}$ be the set of admissible event pairs $(r,j)$ for which
flight $i_r$ can trigger the single type-A check on aircraft $j$.
Use the following A-only variables:
\paragraph{E1-E3. Ordered event variables.}
\begin{itemize}
\item $x_{rj}\in\{0,1\}$: aircraft $j$ operates flight $i_r$;
\item $z_{rj}\in\{0,1\}$: a type-A maintenance event is performed after flight $i_r$ on aircraft $j$;
\item $h_{rj}\ge 0$: cumulative type-A flight time of aircraft $j$ immediately after processing position $r$.
\end{itemize}
Because $r\leftrightarrow i_r$ is one-to-one, the relation to the original
notation is direct: $x_{rj}=x_{i_r j}$ and $z_{rj}=z_{i_r jd}$ for the
corresponding arrival day $d$; the check type is always A in this section.

A maintenance event after flight $i_r$ is assignment-linked and starts at its
arrival time. The corrected airport-continuity, turn-time, and overlap constraints remain
active on the assignment variables. In particular, an assignment is allowed
only when it can be connected to the aircraft's preceding assigned flight and
when the required maintenance duration leaves enough time before the next
assigned departure.

The objective is
\begin{equation}
\min\ \sum_{r\in\mathcal{R}}\sum_{j\in\Pc}c_{rj}x_{rj}
 +\sum_{r\in\mathcal{R}^M}\sum_{j\in\Pc}mc_{rjA}\,z_{rj}.
\label{eq:ordered_event_objective}
\end{equation}

The event formulation numbers its route constraints independently as follows:
E4 is flight coverage, E5 is non-home continuity, E6 is home-airport
continuity and turn-time balance, and E7 is pairwise non-overlap. These are
implemented by \texttt{\_add\_c1\_coverage}, \texttt{\_add\_c23\_turn}, and
\texttt{\_add\_overlap}. E1--E3 are the binary/continuous variable domains defined
above.

The event constraints separate maintenance blocking (E14), event assignment
(E8), and maintenance capacity (E15).

\paragraph{E8. Event assignment.}
Events satisfy
\begin{equation}
z_{rj}\le x_{rj},
\qquad \forall (r,j)\in\mathcal{Z},
\label{eq:ordered_event_assignment}
\end{equation}
The integer event-count link in~\eqref{eq:ordered_c11_recovery} is an auxiliary C11
aggregation, not a separate E family. Timestamp-based incompatibility rows are E14;
they prevent an event from being followed by a departure before the
maintenance duration and turn time have elapsed.


\subsection{Prefix-state propagation}

Let $T_F:=T^A_{\max}$ be the maximum flying time between type-A maintenance events and let $\Phi_j$
be the pre-horizon accumulated flying time of aircraft $j$. At every ordered
position, an unassigned flight carries the state forward, an assigned flight
adds its duration, and an event resets the state after its triggering flight.
Using valid position-specific big-$M$ values $M_{rj}$, the transition is
represented by
\paragraph{E9--E13. Prefix-state propagation.}
\begin{align}
h_{rj}&\ge h_{r-1,j}-M_{rj}x_{rj},\\
h_{rj}&\le h_{r-1,j}+M_{rj}x_{rj},\\
h_{rj}&\ge h_{r-1,j}+p_{r}-M_{rj}(1-x_{rj})-M_{rj}z_{rj},\\
h_{rj}&\le h_{r-1,j}+p_{r}+M_{rj}(1-x_{rj})+M_{rj}z_{rj},\\
h_{rj}&\le T_F(1-z_{rj})
\label{eq:ordered_hour_transition}
\end{align}
The first two rows (E9--E10) carry the state at an unassigned position. The next two
rows (E11--E12) impose accumulation for an assigned flight unless an event is selected.
The final row (E13), together with the nonnegative state domain, gives the exact
post-event state:
$h_{rj}=0$ when $z_{rj}=1$, because the maintenance occurs after flight
$i_r$. The final row together with (E3) bounds the state in all cases and prevents a maintenance event
from repairing an arrival that already exceeds $T_F$. 
The initial state is $h_{0j}=\Phi_j$.

\subsection{Maintenance blocking and capacity constraints}

Let's assume that every maintenance event starts immediately after its triggering flight. Let
 $\delta_A$ be the duration of the type-A check and let $\Delta t$ be the
minimum post-maintenance turn time.

\paragraph{E14. Maintenance blocking (C8).}
For an event $(r,j)$, define the incompatible following-flight set
\begin{equation}
F^{\mathrm{dep}}_{r}=\left\{\ell\in\F:\ \aD{\ell}=\aI{i_r},\ 
  \tI{i_r}<\tD{\ell}<\tI{i_r}+\delta_A+\Delta t\right\}.
\end{equation}
If $z_{rj}=1$, every selected incompatible flight is blocked:
\begin{equation}\label{eq:E14}
x_{\ell j}+z_{rj}\leq1,
\qquad \forall \ell\in F^{\mathrm{dep}}_{r},\ j\in\Pc.
\end{equation}
Only feasible $(\ell,j)$ assignment arcs are generated. Restricting $\ell$ to departures from
the station $\aI{i_r}$ where the check takes place loses nothing: by E5--E6 the aircraft cannot
leave from anywhere else without first operating an intervening flight, which the same family
already blocks.

This constraint also can be rewritten in the next equivalent form:
\begin{equation}\label{eq:E14_rewritten}
\begin{aligned}
M(x_{\ell j}-1) + \tD{\ell} &\ge \tI{i_r} z_{rj}+\delta_A+\Delta t,\\
&\quad \forall \ell\in\F,\ j\in\Pc,\ \tD{\ell}>\tI{i_r}.
\end{aligned}
\end{equation}


\paragraph{E15. Maintenance capacity (C10).}
 Define the active event set
\begin{equation}
\mathcal{E}(a,r)=\left\{(i,j):\aI{i}=a,\quad
\tI{i_r}\leq \tD{i}<\tI{i_r}+\delta_A\right\}.
\end{equation}
At every breakpoint, active maintenance occupancy satisfies
\begin{equation}\label{eq:E15_not_full}
\sum_{(i,j)\in\mathcal{E}(a,r)}z_{i j}
\leq \operatorname{mcap}_{a},
\qquad \forall a\in \MA,\ r\in \mathcal{R}.
\end{equation}

This constraint ensures that the number of simultaneously active maintenance events at any airport 
does not exceed the airport's maintenance capacity for every $r \in \mathcal{R}$.

The constraints \ref{eq:E14} and \ref{eq:E15_not_full} enforce maintenance blocking and capacity limitations, however 
they reduce the solution space by imposing additional restrictions on the feasible assignments of flights and maintenance events.

Potentially, if the plain A instead of going directly to maintenance, waits and during this waiting period 
another plane B arrives and occupies the maintenance slot, 
the constraints will prevent the initial plane from being assigned to maintenance at that time, even if a delayed maintenance could be operationally feasible.


To fix the issue of overly restrictive maintenance assignment due to simultaneous occupancy constraints,
we introduce another variable representing the flexible completion time of maintenance events.
This variable allows the model to adjust the timing of maintenance dynamically,
ensuring that aircraft can be assigned to maintenance even if the initially preferred time slot is occupied.
The new variable is denoted by $s_{rj}\ge0$ and measures the \emph{offset} from the arrival of the
triggering flight $i_r$ to the completion of the check, so that the event occupies the interval
\begin{equation}\label{eq:E16_flex_window}
\bigl[\tI{i_r}+s_{rj}-\delta_A,\ \tI{i_r}+s_{rj}\bigr),
\end{equation}
and $s_{rj}=\delta_A$ recovers the immediate-start behaviour of~\eqref{eq:E14}.

In this case the constraints \ref{eq:E14} and \ref{eq:E15_not_full} turn into the next ones:

\begin{equation}\label{eq:E14_flex}
\begin{aligned}
M(x_{\ell j}-1)+M(z_{rj}-1) + \tD{\ell} &\ge \tI{r}+s_{rj}+\Delta t,\\
\forall \ell\in\F,\ j\in\Pc,\ \aD{\ell}=\aI{i_r}, \ell>r.
\end{aligned}
\end{equation}

This constraint ensures that the completion time of a maintenance event for a given flight and aircraft combination is respected,
taking into account the potential overlap with subsequent flights and the flexibility introduced by the variable $s_{rj}$.
Compared with~\eqref{eq:E14} the blocked window is no longer a constant: an aircraft that defers its
check stays on the ground until the deferred completion, which is the price paid for the extra freedom.
Departures later than $\tI{i_r}+\delta_A+W+\Delta t$ are omitted because $s_{rj}$ cannot reach them.

Additionally, the maintenance capacity constraint \ref{eq:E15_not_full} can be modified to account for the flexible start times.
Writing $T_{rj}=\tI{i_r}+s_{rj}$ for the realised completion time, define the modified event set
\begin{equation}
\mathcal{Q}(a,r,j)=\left\{(i,j'):\aI{i}=a,\quad
s_{rj} - \delta_A < s_{ij'} \le s_{rj}\right\},
\end{equation}
that is, the events at airport $a$ that are in progress at the instant
$s_{rj}-\delta_A$ when event $(r,j)$ begins. Occupancy is then limited by
\begin{equation}\label{eq:E15_flex}
\sum_{(i,j')\in\mathcal{Q}(a,r,j)}\hat{z}_{i j'}
\leq \operatorname{mcap}_{a},
\qquad \forall a\in \MA,\ r\in \mathcal{R},\ j\in\Pc .
\end{equation}
Because every type-A window has the same length $\delta_A$, station occupancy attains its
maximum at the start of some window, so imposing \eqref{eq:E15_flex} at the $|\mathcal{Q}|$
window starts is equivalent to imposing it at every instant. The one-sided membership test is
essential here: counting all events whose windows merely intersect that of $(r,j)$ would
overstate occupancy, since two neighbours of $(r,j)$ need not overlap each other.
Because $\mathcal{Q}(a,r,j)$ depends on the decision variables, \eqref{eq:E15_flex} is
linearised with one membership indicator and one ordering indicator per ordered pair of
potentially conflicting events; the construction is given in
Section~\ref{subsec:flex_complexity}.

We also need the basic condition on $s_{rj}$: the check can neither finish before the triggering
flight has landed and the work has been performed, nor be postponed indefinitely. With
$W$ the maximum admissible deferral,
\begin{equation}\label{eq:E16_flex_s_basic}
\delta_A\,\hat{z}_{rj} \ \le\ s_{rj}\ \le\ (\delta_A+W)\,\hat{z}_{rj},
\qquad \forall r\in \mathcal{R},\ j\in \Pc ,
\end{equation}
and
\begin{equation}\label{eq:E7_flex_s_basic}
\hat{z}_{rj}\ \le\ z_{rj},
\qquad \forall r\in \mathcal{R},\ j\in \Pc .
\end{equation}



The right-hand inequality also forces $s_{rj}=0$ whenever no event is selected, which keeps
\eqref{eq:E14_flex} inactive for non-events.




\subsection{Complete E1--E16 inventory of the implemented A-only model}

The following list gives every event-family component constructed by the current
ordered A-only builder. Legacy C labels are given only as cross-references.

\begin{description}
\item[E1--E3: variable domains.] The variables $x_{rj}$ and $z_{rj}$ are
binary and $h_{rj}$ is nonnegative. These domains correspond to the
assignment, event, and prefix-state variables.

\item[E4: flight coverage (legacy C2).] Each ordered flight is assigned exactly once:
\begin{equation}
\sum_{j\in\Pc}x_{rj}=1,\qquad r\in\mathcal{R}.
\label{eq:ordered_c0_c1}
\end{equation}

\item[E5--E6: aircraft continuity and turn time (legacy C3--C4).] The
corrected equipment-flow balance is applied to $x_{rj}$, including the initial
aircraft position and the minimum turn time. The implementation is
\texttt{e5\_e6\_continuity\_turn}.

\item[E7: pairwise non-overlap (legacy C5).] For every temporally conflicting flight pair
and aircraft, the builder adds
\begin{equation}
x_{ij}+x_{\ell j}\leq1.
\label{eq:ordered_c4_pairwise}
\end{equation}
These rows are constructed by \texttt{e7\_pairwise\_overlap}; in the current
ground-truth formulation they are the corrected C5 family.

The current ground-truth Sections 01--04 do not define a separate clique
constraint number; any clique strengthening is an implementation strengthening
of C5 rather than a new numbered requirement.

\item[E8: event assignment (legacy C6--C7 and C9).]
The maintenance event variables $z_{rj}$ are
binary and are instantiated only for flights arriving at a maintenance-capable
airport. There is no independent daily $y$ variable in this formulation.
Each admissible arrival can trigger at most one type-A event for an aircraft.

\item[E14: maintenance blocking (legacy C8).] If $z_{rj}=1$, every selected flight whose
departure falls before the event's completion plus turn time is blocked:
\begin{equation}
x_{\ell j}+z_{rj}\leq1.
\label{eq:ordered_c8_block}
\end{equation}
Only economically relevant timestamp conflicts are generated.

\item[E15: maintenance capacity (legacy C10).] At every airport timestamp breakpoint, the
active maintenance events satisfy
\begin{equation}
\sum_{(r,j)\in\mathcal{E}(a,\theta)}z_{rj}
\leq \operatorname{mcap}_{a},
\end{equation}
where $\mathcal{I}(a,\theta)$ contains events whose maintenance intervals
cover $r$. This is implemented by the E15 component
\texttt{e15\_maintenance\_capacity}; it is the maintenance-capacity constraint that was not
explicitly displayed earlier.

\item[E14\_flex--E16\_flex: deferrable maintenance timing (optional variant).] Replacing
E14--E15 by \eqref{eq:E14_flex}, \eqref{eq:E15_flex} and \eqref{eq:E16_flex_s_basic} turns the
completion time into a decision. E1--E13 are untouched, so the variant differs from the
immediate-start model only in how an event occupies its station and blocks its aircraft.
It is implemented by \texttt{e14\_flex\_maintenance\_block},
\texttt{e15\_flex\_separation}, \texttt{e15\_flex\_maintenance\_capacity} and
\texttt{e16\_flex\_completion\_window}.

\item[Auxiliary C11 event-count aggregation.] The binary day indicator is intentionally
not retained. Define the nonnegative integer event count $n_{jd}$ by
\begin{equation}
n_{jd}=\sum_{r:\,\lfloor\tI{i_r}/1440\rfloor=d}z_{rj},
\qquad n_{jd}\in\mathbb{Z}_{+}.
\label{eq:ordered_c11_recovery}
\end{equation}
This preserves C11's aggregation information without imposing the removed
one-event-per-day restriction.

\item[Inactive legacy C12.] Calendar spacing is inactive in the A-only model. It is
implemented only by the C/D calendar extension, which is not enabled here.

\item[E9--E13: cumulative flight-hour state (legacy C13--C13f).] The paper's C13 row is
replaced by the auxiliary prefix state $h_{rj}$, which is
nonnegative, carries forward at unassigned positions, accumulates an assigned
flight when $z_{rj}=0$, resets to zero when $z_{rj}=1$, and is bounded by
$T^A_{\max}$. The daily flight-time definition is C13a and the state rows are
C13b--C13f. These are implemented by
\texttt{event\_hour\_state}; they replace the legacy day-pair C13 row rather
than adding another numbered C-constraint.
\end{description}

The event formulation therefore has no additional business rule hidden under
an unnumbered C-constraint. The auxiliary families are the integer C11 event
counts and the linearization rows for the C13 hour
state. The airport occupancy quantity $q_{a,j}$ proposed in the design note is
not introduced as a variable: its role is represented directly by the C10
timestamp-capacity sums, which is smaller and equivalent for the A-only case.

\subsection{Equivalent A-only implementation strengthening}
\label{subsec:event_implementation_strengthening}

The E1--E15 formulation also admits an equivalent reduced implementation.
The daily event count in~\eqref{eq:ordered_c11_recovery} is not used by any
E-constraint or by the objective, so it can be recovered after solution
rather than represented by integer variables and equalities. With only one
maintenance type, event-type exclusivity is implied by the binary domain of
$z_{rj}$ and is omitted.

The prefix state receives its valid upper bound directly as a variable bound.
Consequently, the additional generated row that repeats this bound is
redundant with E13 and is removed. Identical E15 active-event sets are emitted
only once, and a capacity row is omitted whenever its number of binary terms
does not exceed the capacity and is therefore satisfied for every assignment.
Finally, maximal interval-clique rows replace the pairwise-plus-clique E7
implementation; every conflicting pair belongs to a maximal clique, so the
integer feasible set is unchanged.

The baseline implementation remains
\texttt{PaperEventBasedMILPScheduler}. The equivalent reduced implementation
is \texttt{OptimizedPaperEventBasedMILPScheduler}; keeping both classes makes
the effect of implementation choices independently reproducible. The deferrable variant of
Section~\ref{subsec:flex_complexity} is \texttt{FlexiblePaperEventBasedMILPScheduler}, which
derives from the reduced implementation and overrides only the maintenance-timing families.

Table~\ref{tab:ordered_c1_c13} records the correspondence with the numbered
day-indexed requirements. Some legacy rows disappear because their variables
are eliminated; this is a formulation change, not an omitted feasibility
condition.

\begin{table}[t]
\centering\scriptsize
\caption{Coverage of the numbered C1--C13f requirements in the ordered A-only
	formulation.}
\label{tab:ordered_c1_c13}
\begin{tabular}{p{0.10\linewidth} p{0.30\linewidth} p{0.48\linewidth}}
\toprule
Requirement & Day-indexed role & Ordered A-only counterpart \\
\midrule
C1 & Binary assignment domain & Binary domain on $x_{rj}$. \\
C2 & Flight coverage & Coverage equation on $x_{rj}$ (equivalently $x_{i_rj}$). \\
C3--C4 & Airport continuity and turn time & Corrected equipment-flow balance on assignment variables. \\
C5 & Pairwise non-overlap & Pairwise overlap inequalities on $x$. \\
C6--C7 & Maintenance variables $z,y$ & Event indicator $z_{rj}$; the daily indicator $y$ is not a primitive variable. \\
C8 & Maintenance blocks later flights & Timestamp blocking using $\delta_A+\dt$; the flexible variant blocks until the deferred completion. \\
C9 & Maintenance requires assignment & The event-assignment inequality above. \\
C10 & Maintenance capacity & Direct timestamp-overlap capacity rows, equivalent to an airport occupancy state; disjunctive occupancy rows in the flexible variant. \\
C11 & Aggregate events by aircraft and day & Integer event count $n_{jd}=\sum_{r:\,d(i_r)=d}z_{rj}$. \\
C12 & Calendar-time spacing & Inactive for type-A-only maintenance in this paper. \\
C13 & Paper cumulative-flight-time row & Replaced by the exact ordered prefix state. \\
C13a & Daily assigned flight time & Implicit in the ordered flight sequence. \\
C13b--C13f & Exact cumulative-state recursion and bound & Ordered prefix-state rows~\eqref{eq:ordered_hour_transition}. \\
\bottomrule
\end{tabular}
\end{table}

\begin{proposition}[Ordered prefix-state validity]
\label{prop:ordered_event_correct}
Assume the corrected route constraints define a physically feasible integral
assignment, every event is attached to an assigned flight at a
maintenance-capable airport, and the supplied initial state is valid. Then the
ordered prefix-state formulation satisfies the type-A flight-hour limit.
\end{proposition}

\begin{proof}
Induct on the ordered flight positions. An unassigned position carries the
state unchanged. An assigned position adds its flying time unless an event is
selected, in which case the reset rows set the state to zero after the current
flight. The final transition row and the upper state bound prevent
the threshold from being exceeded before a reset. The induction starts from
$\Phi_j$, so every assigned route satisfies the type-A limit.
\end{proof}

\subsection{Why the ordered formulation is the better method}

\subsection{Theoretical complexity and model size}

Let $n_f=|\F|$, $n_p=|\Pc|$, and let $n_m$ denote the number of flights
that arrive at a maintenance-capable airport. In the A-only ordered model,
the assignment variables contribute at most $n_fn_p$ binary variables, while
the maintenance-event variables contribute at most $n_mn_p$ binary variables.
The aircraft-local prefix state contributes at most $n_fn_p$ continuous
variables; with sparse aircraft eligibility, all three counts are smaller and
are determined by the feasible $(i,j)$ pairs. Thus the principal variable
bound is
\[
O(n_fn_p+n_mn_p),
\]
with no additional day-indexed $y_{jd}$ or occupancy-state $q_{a,j}$ family.
The ordered implementation instead stores one state position per feasible
flight--aircraft pair.

The main linear constraint families are $O(n_f)$ for coverage,
$O(n_fn_p)$ for continuity and event assignment, and $O(n_fn_p)$ for the
prefix-state equations. C11 aggregation contributes $O(n_pD)$ equalities for
the event counts, without restricting their values. Pairwise
non-overlap and timestamp blocking depend on the conflict sets and are
$O(|\mathcal{O}|n_p)$ and $O(|\mathcal{B}|)$, respectively, where
$\mathcal{O}$ is the flight-overlap set and $\mathcal{B}$ is the generated set
of event--following-flight blocking pairs. Timestamp-capacity rows are
generated per airport breakpoint, so their number is proportional to the
number of distinct event timestamps rather than to a full time grid. In the
worst case, $|\mathcal{O}|=O(n_f^2)$ and the sparse conflict families can also
be quadratic, but timetable sparsity usually makes their observed counts much
smaller. The theoretical conclusion is therefore precise: the Event-based
model removes the day dimension from the principal variable families and
keeps prefix-state size linear in feasible flight--aircraft pairs; it does not
guarantee fewer total constraints unless the induced conflict and blocking
sets are also sparse.

The ordered formulation combines event-triggered maintenance with a global
flight order and aircraft-local prefix states. It removes the day-pair
accumulation ambiguity of the legacy cumulative-flight-time row, avoids successor-variable
expansion, and keeps the physical route constraints needed for valid tail
assignments. For type-A maintenance, C12's calendar-spacing role is inactive:
the flight-hour limit is enforced by the exact prefix-state recursion. The
computational study in
\Cref{sec:computational} evaluates this formulation against the legacy 
model on matched horizon and fleet-size grids.

\subsection{Complexity and model size of the flexible-start variant}
\label{subsec:flex_complexity}

The deferrable variant keeps E1--E13 verbatim and therefore inherits the variable bound
$O(n_fn_p+n_mn_p)$ of the immediate-start model. On top of it,
\eqref{eq:E16_flex_s_basic} introduces one continuous offset $s_{rj}$ per admissible event,
that is at most $n_mn_p$ additional continuous variables, so the principal bound is unchanged.
The implementation also carries one aggregated offset per maintenance-capable flight; this is
legitimate because E4 assigns each flight to exactly one aircraft, so at most one $s_{r\cdot}$
is nonzero and the aggregate equals the realised offset. The aggregation costs $n_m$ extra
continuous variables and $n_m$ equalities, and it reduces the capacity linearisation below
from event pairs to flight pairs.

The constraint families change in three places. First, the blocking set
$\mathcal{B}$ of \eqref{eq:E14} is generated from the constant window
$[\tI{i_r},\tI{i_r}+\delta_A+\Delta t)$, whereas \eqref{eq:E14_flex} must cover the widened
window $[\tI{i_r},\tI{i_r}+\delta_A+W+\Delta t)$. Its size therefore grows with the deferral
budget $W$, and the rows change from set-packing inequalities to big-M inequalities, which
weakens the linear relaxation. Second, \eqref{eq:E15_flex} is linearised with one membership
indicator $u_{ee'}$ and one ordering indicator per \emph{ordered} pair $(e,e')$ of events at a
maintenance station, where $u_{ee'}=1$ whenever $T_{e'}\in(T_e-\delta_A,T_e]$. Two disjunctive
rows per ordered pair force $u_{ee'}$ up when the membership holds, and one row per event
bounds $\sum_{e'}u_{ee'}\le\operatorname{mcap}_a-1$ whenever $e$ is selected. Writing
$\mathcal{K}_a$ for the set of such ordered pairs, the capacity family is
$O\bigl(\sum_a|\mathcal{K}_a|\bigr)$ rows and the same order of binary variables, against
$O(n_f)$ breakpoint rows and no new binaries in the immediate-start model. Third, a station
whose windows can never exceed its capacity contributes nothing: the builder sweeps the
half-open windows $[\tI{i_r},\tI{i_r}+\delta_A+W)$ and omits the whole family whenever the
peak potential concurrency does not exceed $\operatorname{mcap}_a$, because the inequality is
then satisfied by every assignment.

Two bounds follow. In the worst case $|\mathcal{K}_a|=O(n_m^2)$, so an unbounded deferral
makes the conflict graph complete and the capacity family quadratic; this is precisely why
$W$ is kept finite. In the practical case, only events whose arrival times lie within
$\delta_A+W$ of one another can conflict, so $|\mathcal{K}_a|$ is proportional to $n_m$ times
the number of arrivals inside one deferral window, which is a timetable-dependent constant.
On the generated family of \Cref{sec:computational} the sweep in fact discharges the capacity
family entirely, so with $W$ equal to one day the measured cost reduces to the deferral
machinery alone: an increase of roughly one fifth in variables and one half in rows relative
to the immediate-start implementation.

Finally, the variant is a relaxation of the blocking requirement. Fixing
$s_{rj}=\delta_A z_{rj}$ satisfies \eqref{eq:E16_flex_s_basic} and reduces
\eqref{eq:E14_flex} to \eqref{eq:E14_rewritten}, so every schedule admitted by E14 is admitted
by E14\_flex. For the capacity requirement the two families are not literal rewritings of one
another: \eqref{eq:E15_not_full} tests the fixed breakpoints $\tI{i_r}$, whereas
\eqref{eq:E15_flex} tests the realised window starts, which is the exact occupancy condition.
On any instance where the occupancy sweep certifies that capacity cannot bind, both families
are vacuous and the inclusion is unconditional; this is the case throughout
\Cref{sec:computational}, where the two models return identical optima on every run.
The additional freedom is not free: an aircraft that defers a check is grounded
until the deferred completion, so the relaxation pays for station capacity with aircraft
availability.
